\documentclass[12pt]{article}
\usepackage[utf8]{inputenc}
\usepackage[T1]{fontenc}
\usepackage[margin=1in]{geometry}
\usepackage{setspace}
\usepackage{xcolor}
\usepackage[colorlinks=true, urlcolor=blue!60!black, linkcolor=blue!60!black]{hyperref}
\setlength{\parindent}{2em}
\setlength{\parskip}{1.25em}
\renewcommand{\baselinestretch}{1.0}
\begin{document}
\begin{center}
{\LARGE \textit{INS v. Chadha}, 462 U.S. 919 (1983)\par}
\vspace{0.8em}
{\large Case Brief --- Four-Page Edition\par}
\end{center}
\vspace{1em}

\section{Memory Jogger}

\textit{Chadha} struck down the one-House legislative veto: after delegating authority to the Executive, Congress must use bicameral passage and presentment to alter legal rights and duties outside the Legislative Branch through legislative action (7--2) (\href{https://supreme.justia.com/cases/federal/us/462/919/}{462 U.S. 919}).

\section{Detailed Case Facts}

Jagdish Rai Chadha, an East Indian born in Kenya holding a British passport, entered the United States on a student visa in 1966 and overstayed it. In deportation proceedings he conceded deportability but sought suspension under \S~244(a)(1) of the Immigration and Nationality Act, which let the Attorney General suspend deportation of a person present for seven years, of good moral character, whose deportation would cause ``extreme hardship.'' In 1974 an Immigration Judge granted suspension and reported it to Congress.

Under \S~244(c)(2), either House could disapprove a suspension by resolution, requiring deportation. On December 16, 1975, three days before the window closed, the House passed H.R. Res. 926, without printed text, debate, or a recorded vote, opposing residence for Chadha and five others out of 340 reported suspensions. The resolution was never sent to the Senate or presented to the President.

\section{Procedural History}

The Immigration Judge reopened proceedings and ordered Chadha deported, and the Board of Immigration Appeals affirmed, each holding it could not pass on the statute's constitutionality. Because the INS agreed that the veto was unconstitutional, the Ninth Circuit invited the House and Senate to participate, then held \S~244(c)(2) unconstitutional. The INS appealed and both Houses petitioned for certiorari. After argument and reargument, the Court affirmed on June 23, 1983.

\section{Judicial Votes}

Affirmed, \textbf{7--2}.
\begin{itemize}
\item \textbf{Majority:} Chief Justice Burger, joined by Justices Brennan, Marshall, Blackmun, Stevens, and O'Connor.
\item \textbf{Concurring in the judgment:} Justice Powell.
\item \textbf{Dissents:} Justice White; Justice Rehnquist, joined by Justice White.
\end{itemize}

\section{Holding}

The case was justiciable, Chadha had standing, and \S~244(c)(2) was severable. On the merits, the one-House veto is an exercise of legislative power that fails the bicameralism and presentment requirements of Article I, \S\S~1 and 7, and is therefore unconstitutional.

\section{Analysis of Opinions}

\textbf{Burger (majority).} The INS was ``aggrieved''; the veto was severable under the INA's severability clause, leaving \S~244 a ``fully operative'' report-and-wait scheme; and no political question barred review. On the merits, the proliferation of 295 veto provisions could not save a procedure contrary to the text. Presentment protects the Executive and guards against improvident laws; bicameralism divides legislative power. Action is legislative if it alters legal rights and duties outside the Legislative Branch, and absent the veto Congress could have deported Chadha only by statute: ``Congress must abide by its delegation of authority until that delegation is legislatively altered or revoked.'' The Constitution's explicit unicameral exceptions (impeachment, trials, appointments, treaties) foreclose implied ones. In Footnote 16, Burger treated executive action under a statute as Article II execution bounded by statutory standards and judicial review.

\textbf{Powell (concurring in the judgment).} The sweeping holding would invalidate hundreds of statutes. He would hold that the House assumed a judicial function by deciding that six named individuals failed statutory criteria, exposing them to ``the tyranny of a shifting majority.''

\textbf{White (dissenting).} The legislative veto is an ``indispensable, political invention'' that lets Congress delegate without abdicating; without it, Congress faces a ``Hobson's choice.'' Agencies make binding rules without presentment, the suspension was effectively a legislative proposal requiring the assent of the President and both Houses, and under \textit{Nixon v. Administrator of General Services} the Executive's functions were not impaired.

\textbf{Rehnquist (dissenting).} The veto was inseverable from the suspension power, so Chadha could not obtain relief.

\textbf{Later developments.} \textit{Process Gas} (1983) extended \textit{Chadha} to two-House vetoes of agency rules; \textit{Bowsher v. Synar}, \textit{Alaska Airlines v. Brock}, and \textit{Metropolitan Washington Airports Authority} built on it.

\section{Examples: Future Applications}

\begin{enumerate}
\item \textbf{Same-Side (Invalid): Two-House Disapproval of AI Rules.} A statute lets both Houses block AI safety rules by concurrent resolution, with an inseverability clause. The resolution fails presentment, and the clause means the rulemaking power falls too.
\item \textbf{Same-Side (Invalid): Committee Approval of Land Sales.} Committees must approve federal land sales. Their refusal alters the legal rights of the Secretary and buyer and is even further from Article I, \S~7; the clause is presumptively severable.
\item \textbf{Opposite-Side (Valid): Report-and-Wait with Joint Resolution.} Rules wait 60 days and Congress may block them only by a presented joint resolution, as under the Congressional Review Act.
\item \textbf{Opposite-Side (Valid): Senate Rejection of a Nominee.} Consent to appointments is an express unicameral exception, not lawmaking that bypasses presentment.
\item \textbf{Fence-Sitter: Approval Before Major Rules Take Effect.} Major rules take effect only if both Houses pass and the President signs an approval resolution. Every binding rule passes Article I, \S~7, yet one House's inaction works like the veto struck down here.
\end{enumerate}

\section{Critique}

\textbf{The Footnote 16 asymmetry.} If agency rules and congressional vetoes both alter the legal rights of persons outside the Legislative Branch, \textit{Chadha}'s test struggles to explain why the Executive may make binding rules without bicameralism and presentment while Congress may not reserve a check on them.

\textbf{Severability distortion.} By severing \S~244(c)(2) rather than striking \S~244, the Court arguably gave the Executive broader, unchecked power than Congress ever agreed to delegate, the concern driving Rehnquist's dissent.

\section{Key Quotations}

\begin{quote}
``It emerges clearly that the prescription for legislative action in Art. I, \S\S~1, 7, represents the Framers' decision that the legislative power of the Federal Government be exercised in accord with a single, finely wrought and exhaustively considered, procedure.''
\end{quote}
\noindent --- Chief Justice Burger, Opinion of the Court, 462 U.S. at 951.

\begin{quote}
``In purporting to exercise power defined in Art. I, \S~8, cl. 4, to `establish an uniform Rule of Naturalization,' the House took action that had the purpose and effect of altering the legal rights, duties, and relations of persons, including the Attorney General, Executive Branch officials and Chadha, all outside the Legislative Branch.''
\end{quote}
\noindent --- Chief Justice Burger, Opinion of the Court, 462 U.S. at 952.

\begin{quote}
``Convenience and efficiency are not the primary objectives---or the hallmarks---of democratic government\ldots''
\end{quote}
\noindent --- Chief Justice Burger, Opinion of the Court, 462 U.S. at 944.

\end{document}
